% ============================================================
%  Chapter 12 — Sums of Subsets, Subspaces, and Direct Sums
% ============================================================
\section{The Sum of Subsets, Subspaces, and Direct Sum}
\label{sec:directsum}

Having seen how to intersect subspaces, we now study how to \emph{combine}
them.  The \emph{sum} of two subspaces is the smallest subspace containing
both, and the more refined notion of a \emph{direct sum} ensures that the
decomposition of a vector is unique.

% ---- Sum of subsets ----------------------------------------
\begin{defbox}[Definition --- Sum of Subsets and Subspaces]
\label{def:sum-subsets}
For subsets $S, T\subseteq V$, their \textbf{sum} is
\[
    S + T = \{s+t \mid s\in S,\; t\in T\}.
\]
If $W_1,\ldots,W_k$ are subspaces of $V$, their \textbf{sum} is
\[
    W_1 + W_2 + \cdots + W_k
    = \{w_1+w_2+\cdots+w_k \mid w_i\in W_i\}.
\]
Note that $W_i\subseteq W_1+\cdots+W_k$ for each $i$, and
$W_1\cup\cdots\cup W_k\subseteq W_1+\cdots+W_k$.
\end{defbox}

% ---- Sum is a subspace -------------------------------------
\begin{propbox}[Proposition --- Sum of Subspaces is a Subspace]
\label{prop:sum-subspace}
If $W_1,\ldots,W_k$ are subspaces of $V$, then
$W_1+\cdots+W_k$ is a subspace of $V$.  It is the smallest subspace of $V$
containing every $W_i$.
\end{propbox}

\begin{proofbox}[Proof --- Sum of Subspaces is a Subspace]
$\mathbf{0}=\mathbf{0}+\cdots+\mathbf{0}\in W_1+\cdots+W_k$.  If
$v=\sum w_i$ and $v'=\sum w_i'$ with $w_i,w_i'\in W_i$, then
$v+v'=\sum(w_i+w_i')\in W_1+\cdots+W_k$.  Scalar multiplication is
similar.  Any subspace containing all $W_i$ must contain all sums of their
elements, so the sum is minimal.
\hfill$\blacksquare$
\end{proofbox}

% ---- Span description of sum -------------------------------
\begin{tipbox}[Remark --- Sum via Bases]
If $B_{W_1}=\{u_1,\ldots,u_s\}$ and $B_{W_2}=\{w_1,\ldots,w_t\}$ are bases
for $W_1$ and $W_2$ respectively, then
\[
    W_1 + W_2 = \spn\{u_1,\ldots,u_s,\, w_1,\ldots,w_t\}.
\]
In particular $W_1+W_2 = \spn(B_{W_1}\cup B_{W_2})$.
\end{tipbox}

% ---- Dimension formula -------------------------------------
\begin{thmbox}[Theorem --- Dimension Formula for Sums]
\label{thm:dim-formula}
For subspaces $W_1, W_2$ of a finite-dimensional vector space $V$:
\[
    \dime(W_1+W_2) = \dime W_1 + \dime W_2 - \dime(W_1\cap W_2).
\]
More generally, $\dime\!\left(\sum_{i=1}^k W_i\right)\le\sum_{i=1}^k\dime W_i$.
\end{thmbox}

\begin{proofbox}[Proof --- Dimension Formula for Sums]
See the original lecture notes for a complete proof
(reference: [104016 L-251222]-P14).

\medskip\noindent
\textit{Sketch.}  Let $r=\dime(W_1\cap W_2)$ and choose a basis
$\{c_1,\ldots,c_r\}$ for $W_1\cap W_2$.  Extend it to bases
$\{c_1,\ldots,c_r,u_1,\ldots,u_s\}$ for $W_1$ and
$\{c_1,\ldots,c_r,w_1,\ldots,w_t\}$ for $W_2$.  One verifies that the
combined list $\{c_1,\ldots,c_r,u_1,\ldots,u_s,w_1,\ldots,w_t\}$ is a basis
for $W_1+W_2$, giving $\dime(W_1+W_2)=(r+s)+(r+t)-r=\dime W_1+\dime W_2-r$.
\hfill$\blacksquare$
\end{proofbox}

% ---- Direct sum definition ---------------------------------
\begin{defbox}[Definition --- Direct Sum]
\label{def:direct-sum}
The sum $W_1+\cdots+W_k$ is a \textbf{direct sum}, written
$W_1\oplus W_2\oplus\cdots\oplus W_k$, if the only way to write
$\mathbf{0}=w_1+\cdots+w_k$ with $w_i\in W_i$ is $w_i=\mathbf{0}$ for all
$i$.  Equivalently, every $v\in W_1\oplus\cdots\oplus W_k$ has a
\emph{unique} decomposition $v=w_1+\cdots+w_k$ with $w_i\in W_i$.
\end{defbox}

% ---- Equivalent conditions for direct sum ------------------
\begin{propbox}[Proposition --- Equivalent Conditions for Direct Sum]
\label{prop:direct-sum-equiv}
Let $W_1,\ldots,W_k$ be subspaces of $V$.  The following are equivalent:
\begin{enumerate}[label=(\roman*)]
    \item $W_1+\cdots+W_k$ is a direct sum.
    \item If $B_i$ is a basis for $W_i$, then $B_1\cup\cdots\cup B_k$ is a
          basis for $W_1+\cdots+W_k$.
    \item $\dime\!\left(\displaystyle\bigoplus_{i=1}^k W_i\right)
          = \displaystyle\sum_{i=1}^k \dime W_i$.
\end{enumerate}
\end{propbox}

\begin{proofbox}[Proof --- Equivalent Conditions for Direct Sum]
See the original lecture notes for a complete proof
(reference: [104016 L-251222]-P7 and P10).

\medskip\noindent
\textit{Sketch.}
(i)$\Rightarrow$(ii): The union of the bases spans the sum.  A relation
$\sum_i\sum_j\alpha_{ij}b_{ij}=\mathbf{0}$ (where $b_{ij}\in B_i$) gives
$\sum_i w_i=\mathbf{0}$ with $w_i=\sum_j\alpha_{ij}b_{ij}\in W_i$, so
each $w_i=\mathbf{0}$ (direct sum), hence all $\alpha_{ij}=0$ (l.i.\ of $B_i$).
(ii)$\Rightarrow$(iii): counts the union of bases.
(iii)$\Rightarrow$(i): if $\sum w_i=\mathbf{0}$ with $w_i\in W_i$, extend each
$\{w_i\}\cup B_i$ argument shows each $w_i=\mathbf{0}$.
\hfill$\blacksquare$
\end{proofbox}

% ---- Two-subspace criterion --------------------------------
\begin{propbox}[Proposition --- Direct Sum of Two Subspaces]
\label{prop:direct-sum-two}
$W_1 + W_2$ is a direct sum if and only if $W_1\cap W_2=\{\mathbf{0}\}$.
\end{propbox}

\begin{proofbox}[Proof --- Direct Sum of Two Subspaces]
($\Rightarrow$) If $v\in W_1\cap W_2$, then $v+(-v)=\mathbf{0}$ is a
decomposition with $v\in W_1$ and $-v\in W_2$.  By the direct-sum condition
$v=\mathbf{0}$.

($\Leftarrow$) Suppose $w_1+w_2=\mathbf{0}$ with $w_i\in W_i$.  Then
$w_1=-w_2\in W_1\cap W_2=\{\mathbf{0}\}$, so $w_1=w_2=\mathbf{0}$.
\hfill$\blacksquare$
\end{proofbox}

% ---- Complement --------------------------------------------
\begin{defbox}[Definition --- Complement of a Subspace]
\label{def:complement}
A subspace $W'\leq V$ is a \textbf{complement} of $W\leq V$ if
\[
    V = W \oplus W'.
\]
That is, $V = W+W'$ and $W\cap W'=\{\mathbf{0}\}$.
\end{defbox}

\begin{tipbox}[Remark --- Dimension of a Complement]
If $W'$ is a complement of $W$ in a finite-dimensional $V$, then by the
direct-sum dimension formula:
\[
    \dime_\F W' = \dime_\F V - \dime_\F W.
\]
A complement always exists (extend a basis of $W$ to a basis of $V$; the
added vectors span $W'$), though it is generally not unique.
\end{tipbox}
